
\subsection{Overview}

At a high level, MTSP-LDP follows the \emph{dissimilarity guided publication}
framework, first proposed for the $w$-event CDP setting~\cite{Kellaris2014} and
later extended to the $w$-event LDP setting~\cite{Ren2022}.
The main idea of this framework is that, the server checks at every timestamp
whether it is more beneficial to approximate the current stream statistics with
the last released statistics, than to publish newly computed stream statistics
with necessary noise.
MTSP-LDP further enhances this framework by proposing an \emph{optimal budget
  allocation strategy} to improve estimation accuracy and a \emph{private binary
  tree structure} to enable multi-task streaming queries.
In more detail, MTSP-LDP consists of four mechanisms, i.e., private
dissimilarity estimation, optimal privacy budget allocation, private adaptive
tree publication, and budget-free multi-task streaming query, as illustrated in
\cref{fig:framework}.

\begin{figure}[htp]
  \centering
  \begin{tikzpicture}[>=stealth,thick,font=\small,
  every node/.style={inner sep=1pt,font=\footnotesize},
  group/.style={draw,dashed,rounded corners=1.5mm},
  rec/.style={draw,minimum height=1cm,text width=2.2cm,align=center,fill=blue!10},
  arr/.style={-Implies,double,double distance=2pt},
  point/.style={circle,fill,minimum size=2pt},
  nd/.style={circle,draw,minimum size=1.5mm,inner sep=0pt}
  ]

  % axis
  \coordinate (O);
  \draw[->] (O) ++(-.7,0) -- ++(5,0) node[right] {time};

  % users
  \node (u1) at (0,-.8) {\includegraphics[width=4mm]{figs/user}};
  \node[right=.3 of u1] (u2) {\includegraphics[width=4mm]{figs/user}};
  \node[right=.3 of u2] (u3) {\includegraphics[width=4mm]{figs/user}};
  \node[right=.3 of u3] (dots) {$\dots$};
  \node[right=.3 of dots] (un) {\includegraphics[width=4mm]{figs/user}};

  \node[group,fit={(u1) (un)},label=right:{$n_t$ active users}] (users) {};

  % timestamps
  \node[point,label=above:{$t$}] (t) at (users.north |- O) {};
  \foreach \i in {1,...,8}{
    \node[point,left=.25*\i of t.center] {};
  }

  \draw[->] (t) -- (users.north);

  \node[draw,fill=yellow!20,minimum width=4cm,minimum height=.5cm,
  below=.5 of users] (FO) {Frequency Oracle (e.g., OUE)};

  \foreach \i in {1,...,3}{
    \draw[->] (u\i) -- (u\i |- FO.north);
  }
  \draw[->] (un) -- (un |- FO.north) node[midway,right=2ex] {perturbed $D_t$};
  \node[below=.5 of dots] {$\dots$};

  \node[rec,below left=.5 and .5 of FO] (M1) {$\CM_1^t$\\dissimilarity estimation};
  \node[rec,right=.6 of M1] (M2) {$\CM_2^t$\\privacy bud-get allocation};
  \node[rec,right=.6 of M2] (M3) {$\CM_3^t$\\private tree publication};
  \node[rec,right=.6 of M3] (M4) {$\CM_4^t$\\multi-task stream query};

  \draw[arr] (M1) -- (M2) node[midway,above=3pt] {$\{\!\hdis_t\!\}$};
  \draw[arr] (M2) -- (M3) node[midway,above=3pt] {$\epsilon_{t,2}$};
  \draw[arr] (M3) -- (M4);

  \draw[->] (FO) -- ++(0,-.5) -| (M1.north);

  \draw[->] (FO) -- ++(0,-.5) -| (M3.north);

  % small tree DT
  \begin{scope}[xshift=-1.8cm,yshift=-2.2cm,thin]
    \node[nd] (n31) {};
    \node[nd,right=.1 of n31] (n32) {};
    \node[nd,right=.1 of n32] (n33) {};
    \node[nd,right=.1 of n33] (n34) {};

    \path (n31)--(n32) node[midway,above=.15,nd] (n21) {};
    \path (n33)--(n34) node[midway,above=.15,nd] (n22) {};
    \path (n21)--(n22) node[midway,above=.15,nd] (n1) {};

    \draw (n21) -- (n1) -- (n22);
    \draw (n31) -- (n21) -- (n32);
    \draw (n33) -- (n22) -- (n34);

    \node[left=.1 of n21] {$\hDT_t$};
  \end{scope}

  % OT_t
  \begin{scope}[xshift=4.7cm,yshift=-5.2cm,thin]
    \node[nd] (n31) {};
    \node[nd,right=.1 of n31] (n32) {};
    \node[nd,right=.1 of n32] (n33) {};
    \node[nd,right=.1 of n33] (n34) {};

    \path (n31)--(n32) node[midway,above=.15,nd] (n21) {};
    \path (n33)--(n34) node[midway,above=.15,nd] (n22) {};
    \path (n21)--(n22) node[midway,above=.15,nd] (n1) {};

    \draw (n21) -- (n1) -- (n22);
    \draw (n31) -- (n21) -- (n32);
    \draw (n33) -- (n22) -- (n34);

    \node[above=0 of n1] {$\OT_t$};
  \end{scope}

  % OT_t-1
  \begin{scope}[xshift=3.3cm,yshift=-5.2cm,thin]
    \node[nd] (n31) {};
    \node[nd,right=.1 of n31] (n32) {};
    \node[nd,right=.1 of n32] (n33) {};
    \node[nd,right=.1 of n33] (n34) {};

    \path (n31)--(n32) node[midway,above=.15,nd] (n21) {};
    \path (n33)--(n34) node[midway,above=.15,nd] (n22) {};
    \path (n21)--(n22) node[midway,above=.15,nd] (n1) {};

    \draw (n21) -- (n1) -- (n22);
    \draw (n31) -- (n21) -- (n32);
    \draw (n33) -- (n22) -- (n34);

    \node[above=0 of n1] {$\OT_{t-1}$};
  \end{scope}

  % OT_t-2
  \begin{scope}[xshift=1.9cm,yshift=-5.2cm,thin]
    \node[nd] (n31) {};
    \node[nd,right=.1 of n31] (n32) {};
    \node[nd,right=.1 of n32] (n33) {};
    \node[nd,right=.1 of n33] (n34) {};

    \path (n31)--(n32) node[midway,above=.15,nd] (n21) {};
    \path (n33)--(n34) node[midway,above=.15,nd] (n22) {};
    \path (n21)--(n22) node[midway,above=.15,nd] (n1) {};

    \draw (n21) -- (n1) -- (n22);
    \draw (n31) -- (n21) -- (n32);
    \draw (n33) -- (n22) -- (n34);

    \node[above=0 of n1] {$\OT_{t-2}$};
    \node[left=.6 of n21] {$\dots$};
  \end{scope}

  % OT_t-w
  \begin{scope}[xshift=-.5cm,yshift=-5.2cm,thin]
    \node[nd] (n31) {};
    \node[nd,right=.1 of n31] (n32) {};
    \node[nd,right=.1 of n32] (n33) {};
    \node[nd,right=.1 of n33] (n34) {};

    \path (n31)--(n32) node[midway,above=.15,nd] (n21) {};
    \path (n33)--(n34) node[midway,above=.15,nd] (n22) {};
    \path (n21)--(n22) node[midway,above=.15,nd] (n1) {};

    \draw (n21) -- (n1) -- (n22);
    \draw (n31) -- (n21) -- (n32);
    \draw (n33) -- (n22) -- (n34);

    \node[above=0 of n1] {$\OT_{t-w+1}$};
  \end{scope}


  \coordinate[below=.4 of M1.south west] (A);
  \coordinate[below=.4 of M4.south east] (B);

  \draw[dashed] (A) -- (B);

  \draw[<-] (M1.south) -- (M1 |- A) node[midway,left] {$\OT_{t-1}$};
  \draw[->] (M3.south) -- (M3 |- A) node[midway,right] {$\OT_t$};
  \draw[<-] (M4.south) -- (M4 |- A) node[midway,right] {$\{\OT_t\}$};

  \draw[<-] (M4.130) -- ++(0,.4) node[above, text depth=.25ex] {query};
  \draw[->] (M4.50) -- ++(0,.4) node[above, text depth=.25ex] {answer};

  \node[below=2ex of A,anchor=north west,align=center] {released private \\ binary trees};
\end{tikzpicture}

  \caption{The framework of MTSP-LDP}
  \label{fig:framework}
  \Description{The figure illustrates the workflow of the MTSP-LDP framework at timestamp $t$. At the top, $n_t$ active users send perturbed data to a Frequency Oracle. The process then flows through four connected modules labeled $M_1^t$ through $M_4^t$.}
\end{figure}
\begin{itemize}
\item {\textbf{Private Dissimilarity Estimation $\CM_1^t$}}.
% \bullethdr{Private Dissimilarity Estimation $\CM_1^t$}.
At each time $t$, the server computes the dissimilarity $\dis_t$ between the
statistics of current user data $D_t$ and the last released stream statistics,
and $\dis_t$ will be used to determine whether to approximate with the previous
release or publish with noise.
The main challenge we need to address in $\CM_1^t$ is how to privately compute
$\dis_t$ as user data $D_t$ is not available to the server in LDP, and instead,
we design an unbiased estimator $\hdis_t$ to estimate $\dis_t$ while preserving
user privacy.
\item {\textbf{Optimal Privacy Budget Allocation $\CM_2^t$}}.
% \bullethdr{Optimal Privacy Budget Allocation $\CM_2^t$}.
Given $\hdis_t$, existing methods~\cite{Kellaris2014,Ren2022} directly compare
$\hdis_t$ with the publication error $\err_t$. If $\hdis_t < \err_t$, then they
approximate with the previous release; otherwise, they publish with noise using a
predetermined privacy budget.
However, we notice that this commonly used strategy is often not optimal,
particularly in the $w$-event setting, as it bases its decision solely on the
current timestamp while neglecting the temporal context within the sliding window.
We therefore propose an optimal privacy budget allocation strategy by cleverly leveraging
previously computed dissimilarities, i.e., $\{\hdis_{t'}\}_{t'\leq t}$.
This strategy will determine whether to approximate or to publish, and how to
optimally allocate the privacy budget at time $t$.

\item {\textbf{Private Adaptive Tree Publication $\CM_3^t$}}.
% \bullethdr{Private Adaptive Tree Publication $\CM_3^t$}.
In MTSP-LDP, the server privately collects user data $D_t$ and represents $D_t$
as a private binary tree to both protect user privacy and support
multi-task streaming query.
This private binary tree participates in dissimilarity estimation in
$\CM_1^t$ as well as streaming query in $\CM_4^t$.
We construct the private binary tree in a data-adaptive way in order to
accurately capture the statistics of user data $D_t$.
The output tree $\OT_t$ is then released and will be used by $\CM_4^t$ to run
streaming queries.

\item {\textbf{Budget-Free Multi-task Streaming Query $\CM_4^t$}}.
% \bullethdr{Budget-Free Multi-task Streaming Query $\CM_4^t$}.
Unlike existing research, which typically focuses on ad-hoc queries, e.g., counting, histogram query, or frequency query,
MTSP-LDP is designed to support multi-task streaming queries, including all
queries defined in \cref{ss:queries}, by utilizing the released private binary
trees $\{\OT_{t'}\}_{t'\leq t}$ without incurring  any additional privacy budget.
We theoretically show that MTSP-LDP satisfies $w$-event LDP.
\end{itemize}

These above sub-mechanisms work in tandem to achieve privacy-preserving data analysis in streaming environments, as illustrated in Algorithm~\ref{alg:mtsp-ldp}. Next, we
will introduce them in detail.
\begin{algorithm}[H]
  \caption{MTSP-LDP}
  \label{alg:mtsp-ldp}
  \KwIn{Data stream $\CD=(D_1,D_2,\ldots)$, privacy budget $\epsilon$,
    sliding window size $w$}
  \KwOut{Respond to queries $Q_t^c, Q_t^r$ or $Q_t^e$ at each time $t$}

  \ForEach{$t=1,2,\ldots$}{
    \eIf{$t \leq w$}{
      $\OT_t\gets\FO(D_t, \epsilon/w)$\tcp*{run FO on $D_t$ with budget
        $\epsilon/w$}
    }{
      \tcp{1. private dissimilarity estimation}
      $\hDT_t\gets\FO(D_t, \epsilon/(2w))$\;
      $\hdis_t\gets\text{dissimilarity}(\hDT_t, \OT_{t-1})$\;
      \tcp{2. optimal privacy budget allocation}
      $\epsilon_{t,2}\gets$ run $\CM_3^t(\{\hdis_t\})$\;
      \tcp{3. private tree publication}
      \eIf{$\epsilon_{t,2}>0$}{
        $\hAT_t\gets\texttt{ATC}(\hDT_t)$\tcp*{adaptive tree construction}
      }{
        $\hAT_t\gets\OT_{t-1}$\;
      }
      $\OT_t\gets\texttt{GroupSmooth}(\hAT_t)$\;
    }
  }
  \tcp{4. multi-task processing}
  Use $\{\OT_{t'}\}_{t'\leq t}$ to respond to queries $Q_t^c, Q_t^r$ or $Q_t^e$.
\end{algorithm}

% As the private binary tree is a fundamental data structure used in our framework, in what follows, we first introduce this data structure, and then elaborate on each mechanism.
